\section*{Scriptural Date and Location}

\begin{dossiertable}
\dossierentry{Offertory}{\label{chronology-offertory}}{Ps.~39:2--4 (40), selected clauses}{Davidic prayer; public thanksgiving}
{\chronodate{offertory}{\chronologyannotation{offertory}}}
{Psalm~39 moves from rescue and public testimony to obedience and renewed petition (vv.~2--18). Davidic reception appears in Theodoret, PG~80:1151B--1154C. The critical date is a Psalter boundary, not this psalm's individual date (USCCB, Psalms introduction); the text establishes no exact writing place, first audience, or Davidic life stage.}
\dossierseparator
\dossierentry{Introit}{\label{chronology-introit}}{Ps.~85:1--4 (86), selected clauses}{Davidic prayer addressed to God}
{\chronodate{introit}{\chronologyannotation{introit}}}
{The superscription names David; petition embraces the nations' worship and returns to danger (Ps.~85:1--17). Theodoret reads David and Hezekiah here (PG~80:1553C--1556B). Those received allocations establish no precise writing place or first audience; the critical date bounds the Psalter, not an individual historical occasion.}
\dossierseparator
\dossierentry{Gradual}{\label{chronology-gradual}}{Ps.~91:2--3 (92)}{Sabbath canticle; praise of God}
{\chronodate{gradual}{\chronologyannotation{gradual}}}
{The Sabbath heading frames praise and the contrast between passing wicked prosperity and lasting righteous fruitfulness (Ps.~91:1--16). Theodoret also receives its future rest (PG~80:1615C--1618A). The psalm supplies no named writer's life stage, precise composition place, or first congregation; the critical boundary applies to the Psalter.}
\dossierseparator
\dossierentry{Alleluia}{\label{chronology-alleluia}}{Ps.~94:3 (95), adapted}{Israel's worship of the great King}
{\chronodate{alleluia}{\chronologyannotation{alleluia}}}
{Praise of the Creator and Shepherd leads to Israel's wilderness warning (Ps.~94:1--11; Heb.~3--4). Theodoret's Josiah allocation is his received interpretation (PG~80:1639B--1642A), not an established composition setting. The first congregation and writing place remain unestablished; the critical date is the Psalter boundary. Ps.~46:3 supplies a verbal parallel, not another appointed dossier.}
\dossierseparator
\dossierevententry{Gospel}{\label{chronology-gospel}}{Luke~7:11--16}{Naim's gate; writing place unestablished}
{\chronodate{gospel}{\chronologyannotation{gospel}}}
{Event: Christ meets a widow's funeral procession at Naim's gate (Luke~7:11--12), between the centurion's healing and John's inquiry (7:1--23).}
{Luke's first composition position comes from the Catholic Encyclopedia XIV (1912), ``The New Testament''; the Roman-imprisonment bound is the Biblical Commission's 1912 response, I--IX. No precise composition city, first community, author's life stage, modern settlement identification, or independent modern critical composition horizon is established here.}
\dossierseparator
\dossierevententry{Communion}{\label{chronology-communion}}{John~6:52, final clause (6:51b)}{Capharnaum synagogue; writing place unestablished}
{\chronodate{communion}{\chronologyannotation{communion}}}
{Event: the Bread of Life discourse follows feeding and travel around the Galilean/Tiberias sea; John~6:60 explicitly names the synagogue at Capharnaum.}
{John~6 joins manna, flesh given for life, refusal and Peter's confession. The first composition position is the Catholic Encyclopedia XIV (1912), ``The New Testament''; the second is VIII (1910), ``Gospel of St. John.'' A precise writing place, first community, author's life stage and independent modern critical composition horizon remain unestablished.}
\dossierseparator
\dossierentry{Epistle}{\label{chronology-epistle}}{Gal.~5:25--6:10}{Paul to the Galatian churches}
{\chronodate{epistle}{\chronologyannotation{epistle}}}
{Paul joins grace, faith, freedom and the Spirit to charity (Gal.~1--6). A Lapide (1614), Galatians proemium/argumentum, holds the preferred position. The Catholic Encyclopedia VI (1909), ``Epistle to the Galatians,'' supplies the first, second and fourth alternatives; XI (1911), ``St. Paul,'' the third. Northern/southern Galatia, writing city, modern recipient geography, missionary stage and an independent modern critical date remain unadjudicated.}
\end{dossiertable}
